\documentclass[a4paper,12pt]{article}

% ------------------------------------------------
% Кодировка, язык, математика
% ------------------------------------------------
\usepackage[T2A]{fontenc}
\usepackage[utf8]{inputenc}
\usepackage[russian]{babel}
\usepackage{amsmath,amssymb,mathtools}
\usepackage{icomma}

% ------------------------------------------------
% Шрифт и типографика
% ------------------------------------------------
\usepackage{PTSans}
\renewcommand{\familydefault}{\sfdefault}
\usepackage{microtype}
\usepackage{setspace}
\onehalfspacing
\usepackage{parskip}

% ------------------------------------------------
% Страница
% ------------------------------------------------
\usepackage[
  a4paper,
  left=20mm,
  right=20mm,
  top=24mm,
  bottom=22mm
]{geometry}
\raggedbottom
\setlength{\headheight}{26pt}

% ------------------------------------------------
% Графика, таблицы, списки
% ------------------------------------------------
\usepackage{tikz}
\usetikzlibrary{calc,arrows.meta,positioning,shapes.geometric}
\usepackage{tabularx,booktabs}
\usepackage{enumitem}
\usepackage{xcolor}

% ------------------------------------------------
% Служебные пакеты
% ------------------------------------------------
\usepackage{hyperref}
\usepackage{fancyhdr}

% ------------------------------------------------
% Палитра
% ------------------------------------------------
\definecolor{accent}{HTML}{008080}
\definecolor{darkaccent}{HTML}{004D4D}
\definecolor{textgray}{HTML}{333333}
\definecolor{softgray}{HTML}{F3F5F5}
\definecolor{linegray}{HTML}{B9C3C3}

\hypersetup{
  colorlinks=true,
  linkcolor=darkaccent,
  urlcolor=darkaccent,
  pdftitle={Чек-лист: признаки делимости, простые множители и НОД},
  pdfauthor={Лёвин Артём Александрович}
}

% ------------------------------------------------
% Колонтитулы
% ------------------------------------------------
\pagestyle{fancy}
\fancyhf{}
\fancyhead[L]{\small\textcolor{textgray}{Ксения Васильченко · 6 класс}}
\fancyhead[R]{\small\textcolor{textgray}{Математика}}
\fancyfoot[C]{\small\textcolor{textgray}{\thepage}}
\renewcommand{\headrulewidth}{0.4pt}
\renewcommand{\headrule}{%
  \hbox to\headwidth{%
    \color{linegray}\leaders\hrule height \headrulewidth\hfill
  }%
}

% ------------------------------------------------
% Команды оформления
% ------------------------------------------------
\newcommand{\sectiontitle}[1]{%
  \vspace{0.25em}%
  {\Large\bfseries\color{darkaccent}#1\par}%
  \vspace{0.3em}%
  {\color{linegray}\rule{\linewidth}{0.6pt}}%
  \vspace{0.5em}%
}

\newcommand{\subsectiontitle}[1]{%
  \vspace{0.2em}%
  {\large\bfseries\color{accent}#1\par}%
  \vspace{0.15em}%
}

\newcommand{\important}[1]{\textbf{\textcolor{darkaccent}{#1}}}
\newcommand{\nodop}{\mathop{\text{НОД}}\nolimits}
\newcommand{\lessondate}{17 сентября 2026 года}

\setlist[itemize]{
  leftmargin=1.6em,
  itemsep=0.2em,
  topsep=0.25em
}

\setlist[enumerate]{
  leftmargin=1.9em,
  itemsep=0.4em,
  topsep=0.3em
}

% ------------------------------------------------
% Стили блок-схем
% ------------------------------------------------
\tikzset{
  flowstart/.style={
    ellipse,
    draw=darkaccent,
    very thick,
    align=center,
    minimum width=4.2cm,
    minimum height=1.1cm,
    fill=softgray,
    font=\small
  },
  flowaction/.style={
    rectangle,
    rounded corners=3pt,
    draw=accent,
    thick,
    align=center,
    text width=8.2cm,
    minimum height=1.15cm,
    fill=white,
    font=\small
  },
  flowdecision/.style={
    diamond,
    aspect=2.15,
    draw=darkaccent,
    thick,
    align=center,
    text width=4.7cm,
    inner sep=1pt,
    fill=softgray,
    font=\small
  },
  flowarrow/.style={
    -{Latex[length=3mm,width=2mm]},
    thick,
    draw=textgray
  }
}

\begin{document}

% ================================================================
% Титульный лист
% ================================================================
\begin{titlepage}
\thispagestyle{empty}
\centering

\vspace*{27mm}

{\small\textcolor{accent}{МАТЕМАТИКА · 6 КЛАСС}\par}

\vspace{8mm}

{\Huge\bfseries\color{darkaccent}Чек-лист\par}

\vspace{4mm}

{\LARGE\bfseries
Признаки делимости, разложение на простые множители и НОД
\par}

\vspace{14mm}

{\large Ксения Васильченко\par}

\vspace{3mm}

{\normalsize \lessondate\par}

\vfill

{\normalsize
\textcolor{textgray}{
Лёвин Артём Александрович эксклюзивно для Ксении Васильченко
}
\par}

\vspace{15mm}

\begin{tikzpicture}[remember picture,overlay]
  \draw[
    accent,
    line width=1.6pt
  ]
    ($(current page.south west)+(12mm,17mm)$)
    .. controls
    ($(current page.south west)+(62mm,31mm)$)
    and
    ($(current page.south east)+(-62mm,7mm)$)
    ..
    ($(current page.south east)+(-12mm,18mm)$);
\end{tikzpicture}

\end{titlepage}

% ================================================================
% Основной чек-лист
% ================================================================
\sectiontitle{1. Короткая вычислительная разминка}

Эти приёмы уже знакомы, но они важны для аккуратной работы с дальнейшими задачами.

\subsectiontitle{Десятичные дроби}

\begin{itemize}
  \item
  При сложении и вычитании записывайте числа так, чтобы
  \important{запятые стояли одна под другой}. Недостающие разряды можно дополнять нулями:
  \[
    2{,}30+1{,}72=4{,}02.
  \]

  \item
  При умножении сначала удобно выполнить действие как с натуральными числами, затем вернуть десятичную запятую по общему числу знаков после запятой:
  \[
    1{,}72\cdot2{,}3=3{,}956.
  \]

  \item
  При делении следите за моментом появления запятой в частном:
  \[
    2{,}42:4=0{,}605.
  \]
\end{itemize}

\important{Самопроверка:} оцените порядок результата. Например, сумма чисел около $2$ и $2$ должна быть около $4$.

\subsectiontitle{Обыкновенные и смешанные дроби}

\begin{itemize}
  \item Для сложения и вычитания дробей сначала приводите их к общему знаменателю.
  \item Смешанное число перед умножением или делением удобно переводить в неправильную дробь:
  \[
    2\frac13=\frac73.
  \]
  \item При делении на дробь умножайте на обратную дробь:
  \[
    2\frac13:1\frac25
    =\frac73:\frac75
    =\frac73\cdot\frac57
    =\frac53
    =1\frac23.
  \]
\end{itemize}

\clearpage

\sectiontitle{2. Как находить неизвестный компонент действия}

Формулы ниже помогают превращать более сложное уравнение в последовательность простых шагов.

\renewcommand{\arraystretch}{1.45}
\begin{center}
\begin{tabularx}{0.96\linewidth}{@{}>{\centering\arraybackslash}p{0.23\linewidth}XX@{}}
\toprule
\textbf{Исходное равенство}
&
\textbf{Как найти первый компонент}
&
\textbf{Как найти второй компонент}
\\
\midrule
$a+b=c$
&
$a=c-b$
&
$b=c-a$
\\
$a-b=c$
&
$a=c+b$
&
$b=a-c$
\\
$ab=c$
&
$a=\dfrac{c}{b}$, $b\ne0$
&
$b=\dfrac{c}{a}$, $a\ne0$
\\
$\dfrac{a}{b}=c$, $b\ne0$
&
$a=bc$
&
$b=\dfrac{a}{c}$ при $c\ne0$
\\
\bottomrule
\end{tabularx}
\end{center}

\subsectiontitle{Пример}

Решим уравнение
\[
  \frac{x+1}{3}=10.
\]

Сначала найдём числитель:
\[
  x+1=10\cdot3=30.
\]

Теперь найдём $x$:
\[
  x=30-1=29.
\]

\important{Проверка:}
\[
  \frac{29+1}{3}=\frac{30}{3}=10.
\]

\important{Полезная привычка:} перед каждым действием словами называйте, что именно сейчас неизвестно: слагаемое, уменьшаемое, вычитаемое, множитель, делимое или делитель.

\clearpage

\sectiontitle{3. Признаки делимости}

Признак делимости позволяет быстро понять, делится ли число без остатка, не выполняя полное письменное деление.

\renewcommand{\arraystretch}{1.45}
\begin{center}
\begin{tabularx}{\linewidth}{@{}p{0.12\linewidth}Xp{0.29\linewidth}@{}}
\toprule
\textbf{На}
&
\textbf{Признак делимости}
&
\textbf{Пример}
\\
\midrule
$2$
&
Последняя цифра чётная: $0$, $2$, $4$, $6$ или $8$.
&
$3458$ делится на $2$.
\\
$3$
&
Сумма всех цифр числа делится на $3$.
&
$456$: $4+5+6=15$, значит, число делится на $3$.
\\
$4$
&
Число, составленное из двух последних цифр, делится на $4$.
&
$1004$: последние две цифры дают $04=4$.
\\
$5$
&
Последняя цифра равна $0$ или $5$.
&
$450$ и $455$ делятся на $5$.
\\
$6$
&
Число одновременно делится на $2$ и на $3$.
&
$1224$ делится и на $2$, и на $3$.
\\
\bottomrule
\end{tabularx}
\end{center}

\subsectiontitle{Как соединять признаки}

Для числа $1224$:
\[
  1+2+2+4=9,
\]
поэтому оно делится на $3$. Последняя цифра $4$ чётная, поэтому число делится на $2$. Следовательно, $1224$ делится и на $6$.

Для проверки делимости на $4$ смотрим только на последние две цифры:
\[
  24:4=6,
\]
поэтому $1224$ делится на $4$.

\important{Типичная ошибка:} при признаке делимости на $4$ проверять две последние цифры по отдельности. Нужно рассматривать именно число, которое они образуют.

\clearpage

\sectiontitle{4. Разложение на простые множители}

\important{Простое число} имеет ровно два натуральных делителя: $1$ и само это число. Примеры простых чисел:
\[
  2,\ 3,\ 5,\ 7,\ 11,\ 13,\ldots
\]

Разложить натуральное число на простые множители означает представить его как произведение простых чисел.

\subsectiontitle{Пример: разложим число 54}

Начинаем с самого маленького подходящего простого делителя:
\[
\begin{array}{r|l}
54&2\\
27&3\\
9&3\\
3&3\\
1&
\end{array}
\qquad\Longrightarrow\qquad
54=2\cdot3\cdot3\cdot3=2\cdot3^3.
\]

На каждом шаге проверяем простые делители по порядку: $2$, затем $3$, затем $5$, $7$ и так далее. После деления снова начинаем проверку полученного числа с маленьких простых делителей.

\subsectiontitle{Ещё один пример}

\[
\begin{array}{r|l}
18&2\\
9&3\\
3&3\\
1&
\end{array}
\qquad\Longrightarrow\qquad
18=2\cdot3^2.
\]

\important{Самопроверка:} перемножьте найденные простые множители. Должно получиться исходное число.

% ================================================================
% Блок-схема 1
% ================================================================
\clearpage
\thispagestyle{empty}

\begin{center}
{\Large\bfseries\color{darkaccent}
Алгоритм: разложение числа на простые множители
\par}

\vspace{10mm}

\begin{tikzpicture}[node distance=10mm]

  \node[flowstart] (start)
  {Начало: дано натуральное число $n>1$};

  \node[flowdecision,below=14mm of start,text width=4.0cm] (done)
  {$n=1$?};

  \node[flowaction,below=16mm of done,text width=7.2cm] (choose)
  {Найдите самый маленький простой делитель $p$, на который $n$ делится без остатка};

  \node[flowaction,below=of choose,text width=7.2cm] (divide)
  {Запишите множитель $p$ и замените $n$ на $n:p$};

  \node[flowstart,right=13mm of done,minimum width=3.7cm] (finish)
  {Конец: разложение\\получено};

  \draw[flowarrow] (start) -- (done);

  \draw[flowarrow]
  (done.east) -- node[above,font=\small]{да} (finish.west);

  \draw[flowarrow]
  (done) -- node[right,font=\small]{нет} (choose);

  \draw[flowarrow] (choose) -- (divide);

  \draw[flowarrow]
  (divide.west) -- ++(-1.6,0) |- (done.west);

\end{tikzpicture}
\end{center}

\clearpage

% ================================================================
% НОД
% ================================================================
\sectiontitle{5. Наибольший общий делитель}

\important{Общий делитель} двух натуральных чисел --- число, на которое каждое из них делится без остатка. \important{Наибольший общий делитель (НОД)} --- самый большой из общих делителей:
\[
  \nodop(a,b).
\]

\subsectiontitle{Как найти НОД через простые множители}

\begin{enumerate}[itemsep=0.1em,topsep=0.15em]
  \item Разложите оба числа на простые множители.
  \item Выберите общие простые множители, учитывая каждое повторение только тогда, когда оно есть в обоих разложениях.
  \item Перемножьте выбранные множители.
\end{enumerate}

Иными словами, каждый общий простой множитель входит в НОД столько раз, сколько он встречается \important{в меньшем количестве} из двух разложений.

\subsectiontitle{Пример 1: $\nodop(18,54)$}

\[
  18=2\cdot3^2,
  \qquad
  54=2\cdot3^3,
  \qquad
  \nodop(18,54)=2\cdot3^2=18.
\]

\subsectiontitle{Пример 2: $\nodop(12,18)$}

\[
  12=2^2\cdot3,
  \qquad
  18=2\cdot3^2,
  \qquad
  \nodop(12,18)=2\cdot3=6.
\]

Во втором примере у числа $12$ две двойки, а у числа $18$ только одна. Поэтому в НОД входит только одна двойка. Аналогично с тройками.

\important{Быстрая проверка:} найденный НОД должен делить оба числа без остатка и не может быть больше меньшего из них.

% ================================================================
% Блок-схема 2
% ================================================================
\clearpage
\thispagestyle{empty}

\begin{center}
{\Large\bfseries\color{darkaccent}
Алгоритм: поиск НОД через простые множители
\par}

\vspace{9mm}

\begin{tikzpicture}[node distance=9mm]

  \node[flowstart] (start)
  {Начало: даны натуральные числа $a$ и $b$};

  \node[flowaction,below=of start,text width=7.2cm] (factor)
  {Разложите $a$ и $b$ на простые множители};

  \node[flowaction,below=of factor,text width=7.2cm] (init)
  {Начните произведение для НОД с числа $1$};

  \node[flowdecision,below=14mm of init,text width=4.2cm] (common)
  {Есть общий неиспользованный простой множитель?};

  \node[flowaction,below=17mm of common,text width=7.2cm] (take)
  {Умножьте НОД на этот множитель и отметьте по одному его вхождению в каждом разложении};

  \node[flowstart,right=12mm of common,minimum width=3.8cm] (finish)
  {Конец:\\получен $\nodop(a,b)$};

  \draw[flowarrow] (start) -- (factor);
  \draw[flowarrow] (factor) -- (init);
  \draw[flowarrow] (init) -- (common);

  \draw[flowarrow]
  (common) -- node[right,font=\small]{да} (take);

  \draw[flowarrow]
  (take.west) -- ++(-1.7,0) |- (common.west);

  \draw[flowarrow]
  (common.east) -- node[above,font=\small]{нет} (finish.west);

\end{tikzpicture}
\end{center}

\clearpage

\sectiontitle{6. НОД и сокращение дробей}

Если числитель и знаменатель дроби имеют общий делитель, дробь можно сократить, разделив числитель и знаменатель на одно и то же число. Деление сразу на НОД даёт несократимую дробь за один шаг.

\subsectiontitle{Пример: сократим $\dfrac{324}{432}$}

Разложим числа:
\[
  324=2^2\cdot3^4,
  \qquad
  432=2^4\cdot3^3.
\]

Берём общие множители в меньших степенях:
\[
  \nodop(324,432)=2^2\cdot3^3=108.
\]

Теперь сокращаем:
\[
  \frac{324}{432}
  =\frac{324:108}{432:108}
  =\frac34.
\]

\subsectiontitle{Полезный способ экономить вычисления}

Если при сравнении разложений Вы отмечаете все множители, вошедшие в НОД, то произведение оставшихся множителей показывает результат деления исходного числа на НОД.

В примере выше после удаления общих множителей у $324$ остаётся $3$, а у $432$ остаётся $2^2=4$. Поэтому сразу видно:
\[
  \frac{324}{432}=\frac34.
\]

\subsectiontitle{Типичные ошибки}

\begin{itemize}
  \item В НОД включён множитель, который встречается только в одном разложении.
  \item Общий множитель взят больше раз, чем он встречается хотя бы в одном из чисел.
  \item При разложении допущена ошибка в промежуточном делении.
  \item При сокращении числитель и знаменатель поделены на разные числа.
  \item После вычисления НОД не выполнена простая проверка делимости обоих исходных чисел.
\end{itemize}

% ================================================================
% Тренировочный блок
% ================================================================
\clearpage
\sectiontitle{7. Тренировочный блок}

Решайте задания по порядку. В заданиях на НОД показывайте разложение чисел на простые множители.

\begin{enumerate}

  \item
  \textbf{Вычислительная разминка.} Найдите значения:
  \begin{itemize}[leftmargin=2.2em]
    \item[а)] $3{,}40+1{,}65$;
    \item[б)] $1{,}25\cdot2{,}4$;
    \item[в)] $3{,}24:4$;
    \item[г)] $2\dfrac14:1\dfrac12$.
  \end{itemize}

  \item
  \textbf{Признаки делимости.}
  \begin{itemize}[leftmargin=2.2em]
    \item[а)] Определите, на какие из чисел $2$, $3$, $4$, $5$, $6$ делится число $1224$.
    \item[б)] В числе $12\square4$ вместо квадрата поставьте все цифры, при которых число одновременно делится на $3$ и на $4$.
  \end{itemize}

  \item
  \textbf{Разложение на простые множители.} Разложите числа $72$ и $90$ на простые множители.

  \item
  \textbf{НОД.} Найдите
  \[
    \nodop(48,72)
  \]
  через разложение на простые множители.

  \item
  \textbf{НОД и сокращение дроби.} Найдите
  \[
    \nodop(414,504),
  \]
  затем сократите дробь
  \[
    \frac{414}{504}.
  \]

\end{enumerate}

% ================================================================
% Ответы
% ================================================================
\clearpage
\sectiontitle{8. Ответы}

\renewcommand{\arraystretch}{1.5}
\begin{center}
\begin{tabularx}{\linewidth}{@{}p{0.10\linewidth}X@{}}
\toprule
\textbf{№}
&
\textbf{Ответ}
\\
\midrule
1
&
а) $5{,}05$; б) $3$; в) $0{,}81$; г) $\dfrac32=1\dfrac12$.
\\
2
&
а) $1224$ делится на $2$, $3$, $4$ и $6$; на $5$ не делится.\newline
б) $2$ и $8$: числа $1224$ и $1284$ одновременно делятся на $3$ и на $4$.
\\
3
&
$72=2^3\cdot3^2$; \quad $90=2\cdot3^2\cdot5$.
\\
4
&
$48=2^4\cdot3$, $72=2^3\cdot3^2$, поэтому $\nodop(48,72)=2^3\cdot3=24$.
\\
5
&
$414=2\cdot3^2\cdot23$, $504=2^3\cdot3^2\cdot7$, поэтому $\nodop(414,504)=18$ и
$\dfrac{414}{504}=\dfrac{23}{28}$.
\\
\bottomrule
\end{tabularx}
\end{center}

\vfill

{\small
\textcolor{textgray}{
Контрольный ориентир: Вы умеете объяснить своими словами признак делимости, разложить число на простые множители, выбрать общие множители для НОД и проверить полученный ответ делением.
}
\par}

\end{document}